\chapter{Introduction}
\label{chap:introduction}

% ~1100-1300 words in total (reference thesis: 1131).

\section{Motivation}
\label{sec:motivation}
% Knowledge graphs are not static: YAGO, DBpedia and Wikidata publish new versions
% regularly, and each release adds, removes and rewrites facts. Consumers of a KG
% (search, QA, recommendation) inherit those changes blindly.
% Argue: we have plenty of metrics that describe ONE snapshot, and almost nothing
% that describes the DIFFERENCE between two, at a scale where the graphs are
% billions of triples.
% TODO

\section{Research Objectives}
\label{sec:research_objectives}
% State them as answerable questions, e.g.
%   RQ1  Which metrics characterise a KG snapshot, and can they be computed
%        efficiently over a billion-triple graph through a SPARQL endpoint?
%   RQ2  How do those metrics change between two versions of the same KG, and
%        where in the graph does the change concentrate?
%   RQ3  Can the comparison of two versions be made efficient enough to run on
%        commodity hardware rather than a cluster?
%   RQ4  What do the same metrics say when applied across two different KGs?
% TODO

\section{Contributions}
\label{sec:contributions}
% Be precise about what is yours versus what is adopted:
%   - a catalogue of 13 metrics over three analysis levels, implemented as
%     SPARQL + Rust and runnable against any QLever endpoint;
%   - entropy-weighted PageRank, a variant in which every edge is weighted by the
%     property entropy of its predicate (Chapter 3);
%   - an integer-encoded triple diff with a subject-window scoping scheme, which
%     makes version comparison tractable without a global sort;
%   - a dashboard that puts metric values of several versions side by side;
%   - an evaluation over YAGO and DBpedia.
% Adopted, NOT claimed: ETImp and classic PageRank (Knowgly baselines).
% TODO

\section{Organization}
\label{sec:organization}
% One short paragraph naming each remaining chapter.
% TODO
